\section{Related work}\label{sec:related}

\subsection{Calibration of computer models with a discrepancy}

Kennedy and O'Hagan~\cite{ko} framed the calibration of a computer model as Bayesian inference on its parameters with an
additive model discrepancy, a Gaussian process in the inputs, so that predictions carry the uncertainty of the parameters and
of the model's error. Higdon et al.~\cite{higdon} added an emulator for expensive simulators, and Bayarri et al.~\cite{bayarri}
built a validation framework around the same model, with inferential statements about the predictive error in untested
conditions. The framework has a known weakness: the parameters and the discrepancy are confounded, so the posterior of a
parameter depends on the discrepancy prior. Brynjarsd\'ottir and O'Hagan~\cite{brynj} showed that ignoring the discrepancy
biases the parameters and the predictions, and that with it the parameters are learned only through meaningful priors, while
predictions inside the range of the data can still be good. Arendt et al.~\cite{arendt} explained when the two can be told
apart and when they cannot, and in a companion paper~\cite{arendt2} showed that several responses that share a parameter
improve its identifiability by an amount that depends on the responses. Wong et al.~\cite{wong} gave an identifiable
frequentist parameterization with a two-step estimator, and Sung and Tuo~\cite{sungtuo} reviewed the framework and the
experimental designs it calls for. The earlier CrusherBayes paper~\cite{cb002} made the confounding exact for a linear
discrepancy: the information about a scalar is cut in the directions of the discrepancy, and on the HP500 lane the capacity
base is confounded with the throughput discrepancy at every value of the scalars.

\subsection{What the calibration parameter should mean}

A second line of work asks which value of the parameter calibration should target when the model is wrong. Tuo and
Wu~\cite{tuowu} defined it as the value that minimizes the $L_2$ distance between the model and the true process, proved that
the Kennedy and O'Hagan estimator can miss it, and gave an $L_2$ calibration that is asymptotically normal and semiparametric
efficient. Tuo~\cite{tuo} showed that the inconsistency is removed by projecting the model's gradient out of the discrepancy's
kernel; Plumlee~\cite{plumlee} reached the same end with a discrepancy prior orthogonal to the gradient of the computer
model, which narrows a posterior that the standard prior leaves suboptimally broad; and Plumlee and Joseph~\cite{plumleejoseph}
gave the orthogonal Gaussian process in closed form. Xie and Xu~\cite{xiexu} proved a Bernstein-von Mises theorem for a
Bayesian projected calibration. Gu and Wang~\cite{guwang} took a different route with the scaled Gaussian stochastic process,
which penalizes the mean square of the discrepancy so that the calibrated model alone predicts well. Woody et al.~\cite{woody}
argued that for extrapolative prediction no inference on the discrepancy is needed and used a Gibbs posterior with a tuned
loss. These works concern the parameter: its consistency, its efficiency, its meaning. None compares the risk of the
calibrated predictor with that of an empirical regression in the same inputs, which is the comparison a practitioner with ten
surveys has to make.

\subsection{Ridge regression, misspecification and model selection}

The predictive mean of a calibration with a linear discrepancy is a ridge regression~\cite{hoerl} toward a target set by the
model~\cite{cb002}. The risk of ridge regression is classical in the fixed design~\cite{hastie}, and Hsu et al.~\cite{hsu} gave
random-design bounds for least squares and ridge with an approximation-error term that plays the role of the plant's curvature
here. The comparison of a regression on two regressors with one on three is Mallows's~\cite{mallows}: a regressor pays when
its signal exceeds its variance. Proposition~\ref{prop:line} combines these with the engine's offset, which none of them has.
For comparing two predictors on the same data, Diebold and Mariano~\cite{diebold} gave tests of equal predictive accuracy; in
a simulation the replicates are independent and the paired difference of errors has a plain Monte Carlo standard error. What
these results do not cover is a shrinkage target that is itself a fitted physical model with a curvature of its own: ridge
theory treats the target as a fixed vector, and the misspecification bounds treat the approximation error of the regression as
a nuisance, not as something a second model can partly predict. Theorem~\ref{thm:risk} adds that term.

\subsection{Physics-informed against data-driven predictors}

The closest empirical study we found is that of Hotvedt et al.~\cite{hotvedt}, who tested by simulation of a virtual flow
meter whether a grey-box model beats a physics-based one under model mismatch and a data-driven one with little data. They
gave no risk formula and no calibration prior. In comminution, Yamashita et al.~\cite{yamashita} classified 61 works on cone
crushers as population-balance, empirical or data-driven; Duarte et al.~\cite{duarte} calibrated a Whiten cone-crusher model
by least squares and validated it on a data set left out of the fit, and found that the most recent calibrations of that model
date from the 1990s; Rocha et al.~\cite{rocha} fitted the HP500 model the engine uses, with a capacity that is a power law
in the setting. A discrete-element simulation could serve as a known truth, but one operating point of a quarter section of a
secondary cone crusher took 172 to 404 CPU-hours in the thesis of Quist~\cite{quistthesis}, whose quarter-section setting
Quist and Evertsson~\cite{quist} published, which rules out the hundreds of operating points a risk study needs.

\subsection{Design for calibration and the checking of samplers}

Bayesian experimental design chooses the design that maximizes an expected utility; Chaloner and Verdinelli~\cite{chaloner}
unified the linear case, Ryan et al.~\cite{ryan} reviewed the computational algorithms, Huan and Marzouk~\cite{huan} gave a
simulation-based method for nonlinear models, and Rainforth et al.~\cite{rainforth} framed the expected information gain as
the mutual information between parameters and data. Raue et al.~\cite{raue} separated structural from practical
non-identifiability with the profile likelihood, the distinction Proposition~\ref{prop:priors} uses. To check that a sampler
returns the posterior of its model, Talts et al.~\cite{talts} proposed simulation-based calibration, in which the rank of a
prior draw among posterior draws fitted to data simulated from it must be uniform; Modr\'ak et al.~\cite{modrak} showed that
ranks of the parameters alone cannot detect a posterior equal to the prior and that the likelihood of the data is a test
quantity that can. Gelman et al.~\cite{gelman} placed fake-data simulation in the Bayesian workflow, and Morris et
al.~\cite{morris} set the standard of planning and reporting a simulation study that Section~\ref{sec:design} follows.

\subsection{Where this paper stands}

The search behind this paper (web and bibliographic databases, 2026-11-19) found no derivation of when a calibrated physics
model beats an empirical regression in the same regressors as a function of the plant's curvature over the operating range,
no evaluation of a mineral-processing calibration on virtual plants with a known truth, and no Bayesian calibration of a
Whiten cone-crusher model other than~\cite{cb002}. The closest results, those on the calibration parameter and those on ridge
risk, each lack the other half: the former do not compare predictors, the latter have no engine. This paper joins them
for a discrepancy linear in a few regressors, where every quantity is exact. The restriction is the one the released method
makes, and it is what lets the risk be written without asymptotics: with a Gaussian-process discrepancy the plane would be
replaced by a kernel smoother, and the curvature gain of Corollary~\ref{cor:align} would compete with the kernel's own ability
to follow the plant, a comparison that the rates of Tuo~\cite{tuo} and of Xie and Xu~\cite{xiexu} describe only as the number
of surveys grows.
